\documentclass[11pt,a4paper]{article}

% ---------- Packages ----------
\usepackage[utf8]{inputenc}
\usepackage[T1]{fontenc}
\usepackage{lmodern}
\usepackage[margin=2.3cm]{geometry}
\usepackage{amsmath,amssymb}
\usepackage{booktabs}
\usepackage{array}
\usepackage{xcolor}
\usepackage{titlesec}
\usepackage{fancyhdr}
\usepackage{enumitem}
\usepackage{hyperref}
\usepackage{graphicx}
\usepackage{tabularx}
\usepackage{ragged2e}
\usepackage{pgfplots}
\usepackage{listings}
\usepackage{tcolorbox}
\usepackage{longtable}
\usepackage{siunitx}
\pgfplotsset{compat=1.18}
\tcbuselibrary{listings,skins}
\sisetup{group-separator={,},group-minimum-digits=4}

% ---------- Colors ----------
\definecolor{gcibblue}{RGB}{12,45,90}
\definecolor{gcibgold}{RGB}{180,140,30}
\definecolor{ghred}{RGB}{140,20,20}
\definecolor{ghgold}{RGB}{200,160,40}
\definecolor{ghgreen}{RGB}{0,100,50}
\definecolor{codebg}{RGB}{247,247,250}
\definecolor{codekw}{RGB}{0,90,160}
\definecolor{codestr}{RGB}{160,60,20}
\definecolor{codecom}{RGB}{90,120,90}

% ---------- Author metadata ----------
\newcommand{\authorname}{Daniel Agyekum Amakye}
\newcommand{\authortitle}{Quantitative Analyst --- Risk Management}
\newcommand{\authordesk}{Quantitative Risk \& Financial Engineering}
\newcommand{\authoremail}{janetobosuayaa@gmail.com}
\newcommand{\reporef}{github.com/DanielAgyekumAmakyeGh/ghana-fx-var}

% ---------- Hyperref ----------
\hypersetup{
    colorlinks=true, linkcolor=gcibblue,
    urlcolor=gcibblue, citecolor=gcibblue,
    pdftitle={Ghana FX VaR: Illustrative vs. Real-Data Calibration},
    pdfauthor={\authorname},
    pdfsubject={USD/GHS Value-at-Risk under Basel III / Bank of Ghana framework}
}

% ---------- Section formatting ----------
\titleformat{\section}
  {\normalfont\Large\bfseries\color{gcibblue}}{\thesection}{1em}{}
\titleformat{\subsection}
  {\normalfont\large\bfseries\color{gcibblue}}{\thesubsection}{1em}{}

% ---------- Header / Footer ----------
\pagestyle{fancy}
\fancyhf{}
\fancyhead[L]{\small\color{gcibblue}\textbf{GCIB} \;|\; \authordesk}
\fancyhead[R]{\small\color{gcibblue}Ghana Market Risk Report}
\fancyfoot[L]{\small\color{gcibblue}\authorname}
\fancyfoot[C]{\small\thepage}
\fancyfoot[R]{\small\color{gcibblue}\authortitle}
\renewcommand{\headrulewidth}{0.4pt}
\renewcommand{\footrulewidth}{0.4pt}

% ---------- Custom commands ----------
\newcommand{\ghs}[1]{\textnormal{GH₵\,#1}}
\newcolumntype{L}[1]{>{\RaggedRight\arraybackslash}p{#1}}

% ---------- Python listing style ----------
\lstdefinestyle{python}{
    language=Python,
    backgroundcolor=\color{codebg},
    basicstyle=\ttfamily\footnotesize,
    keywordstyle=\color{codekw}\bfseries,
    stringstyle=\color{codestr},
    commentstyle=\color{codecom}\itshape,
    numbers=left, numberstyle=\tiny\color{gray}, numbersep=8pt,
    frame=single, rulecolor=\color{gcibblue!40},
    showstringspaces=false, breaklines=true, tabsize=4, captionpos=b
}

% ============================================================
\begin{document}

% ---------- Title Page ----------
\begin{titlepage}
    \centering
    \vspace*{1.2cm}
    {\color{gcibblue}\rule{\textwidth}{1.5pt}}\\[0.4cm]
    {\Huge\bfseries\color{gcibblue} Ghana Commercial \& Investment Bank}\\[0.3cm]
    {\Large\color{gcibgold} Quantitative Risk \& Financial Engineering}\\[0.2cm]
    {\color{gcibblue}\rule{\textwidth}{1.5pt}}\\[1.2cm]

    {\Huge\bfseries Ghana Market Risk Report}\\[0.5cm]
    {\LARGE FX Value-at-Risk: Illustrative vs.\ Real-Data Calibration}\\[0.3cm]
    {\large USD/GHS Position --- Basel III / Bank of Ghana Framework}\\[1.4cm]

    \begin{tcolorbox}[
        width=0.82\textwidth,
        colback=codebg,
        colframe=gcibblue,
        boxrule=1pt, arc=3pt,
        left=14pt, right=14pt, top=10pt, bottom=10pt
    ]
        \centering
        {\large\bfseries\color{gcibblue} Prepared by}\\[0.35cm]
        {\Large\bfseries \authorname}\\[0.25cm]
        {\color{gcibgold}\rule{6cm}{0.6pt}}\\[0.25cm]
        {\large \authortitle}\\[0.15cm]
        {\itshape \authordesk}\\[0.35cm]
        {\small \authoremail}
    \end{tcolorbox}

    \vspace{1.0cm}

    \begin{tabular}{rl}
        \textbf{Desk:}          & Market Risk \\[0.2cm]
        \textbf{Requested by:}  & Treasury / ALM \\[0.2cm]
        \textbf{Regulator:}     & Bank of Ghana (BoG) \\[0.2cm]
        \textbf{Classification:} & Internal --- Restricted \\[0.2cm]
        \textbf{Report Ref:}    & GCIB/MR/FX/2025/002 \\[0.2cm]
        \textbf{Repository:}    & \texttt{\reporef} \\[0.2cm]
        \textbf{Date:}          & \today \\
    \end{tabular}

    \vfill
    {\small\itshape Prepared under the Bank of Ghana's Risk Management Directive and Basel II/III Capital Adequacy framework.}
\end{titlepage}

% ---------- Executive Summary ----------
\section{Executive Summary}

This report presents a dual-calibration FX Value-at-Risk (VaR) analysis
for a USD~20 million position at GCIB. Two complementary engines are
developed:

\begin{enumerate}[leftmargin=1.6em]
    \item An \textbf{illustrative engine} using a controlled parameter
          set (spot 12.50, volatility 1.20\%/day) to demonstrate the
          mechanics of parametric VaR, Expected Shortfall, Kupiec
          backtesting, and portfolio VaR.
    \item A \textbf{real-data engine} calibrated against 2{,}725 daily
          observations of the Bank of Ghana USD/GHS reference rate
          (2015--2025) to test whether the framework survives real
          Ghanaian market conditions.
\end{enumerate}

\noindent\textbf{Headline finding.} The parametric VaR model \emph{fails
the Kupiec Proportion-of-Failures test outright} on real data: 60 breaches
against 24.75 expected over 2{,}475 rolling 250-day windows
(\(LR = 36.27\) against a \(\chi^2_{1,0.95}\) critical value of 3.84).
The parametric assumption understates Expected Shortfall by a factor of
2.14\(\times\) relative to the empirical tail. The bank is fictional; the
methodology and the failure mode are real.

% ---------- Section 2 ----------
\section{Problem Statement}

\noindent\textbf{Requesting unit:} Treasury / ALM

\begin{quote}
\itshape
``We have a USD 20 million FX position. What is our potential loss tomorrow?''
\end{quote}

\noindent The USD/GHS rate is one of the most volatile emerging-market
pairs in Sub-Saharan Africa. Between 2015 and 2025, the cedi ranged from
3.20 to 16.42 against the dollar, with annualised volatility varying
across regimes from 0.88\% (2021) to 30.79\% (2022) --- a 35\(\times\)
spread. A robust VaR framework must both \emph{measure} this risk and be
\emph{validated} against real data.

% ---------- Section 3 ----------
\section{Methodology}

Both engines share the same core methodology. Only the calibration differs.

\subsection{Parametric VaR}
\begin{equation}
\text{VaR}_{1d} = V \times \sigma \times z_{\alpha}
\end{equation}

\subsection{Expected Shortfall (normal case)}
\begin{equation}
\text{ES}_{1d} = V \times \sigma \times \frac{\phi(z_{\alpha})}{1-\alpha}
\end{equation}

\subsection{Historical Simulation VaR}
\begin{equation}
\text{VaR}^{\text{hist}} = \text{quantile}_{1-\alpha}\!\left(-r_t \cdot V\right)
\end{equation}

\subsection{GARCH(1,1)}
\begin{equation}
\sigma_t^2 = \omega + \alpha_1 \varepsilon_{t-1}^2 + \beta_1 \sigma_{t-1}^2
\end{equation}
with Student-\(t\) errors to capture fat tails.

\subsection{Kupiec Proportion-of-Failures}
\begin{equation}
LR_{POF} = -2 \ln\!\left[
\frac{(1-p)^{n-x} p^{x}}
     {\left(1-\tfrac{x}{n}\right)^{n-x} \left(\tfrac{x}{n}\right)^{x}}
\right] \sim \chi^2_1
\end{equation}
where \(p = 1-\alpha\) is the expected breach rate and \(x\) the observed count.

\subsection{Portfolio VaR}
\begin{equation}
\sigma_P = \sqrt{w^\top \Sigma\, w}, \qquad
\text{VaR}_P = z_\alpha \times \sigma_P
\end{equation}

% ---------- Section 4 ----------
\section{Calibration 1: Illustrative Engine}

\subsection{Inputs}

\begin{table}[h!]
\centering
\renewcommand{\arraystretch}{1.3}
\begin{tabular}{L{7cm} L{7cm}}
\toprule
\textbf{Input} & \textbf{Value} \\
\midrule
Currency pair              & USD/GHS \\
Spot rate                  & 12.50 \\
Position (base currency)   & USD 20{,}000{,}000 \\
Position value             & \ghs{250{,}000{,}000} \\
Assumed daily volatility   & 1.20\% \\
Confidence level           & 99\% \\
Horizon                    & 1 day (Basel: 10 days) \\
\bottomrule
\end{tabular}
\end{table}

\subsection{Results}

\begin{table}[h!]
\centering
\renewcommand{\arraystretch}{1.3}
\begin{tabular}{L{7cm} r}
\toprule
\textbf{Metric} & \textbf{Value (GHS)} \\
\midrule
Parametric VaR (99\%, 1-day)     & 6{,}979{,}044 \\
Parametric VaR (99\%, 10-day)    & 22{,}069{,}674 \\
Expected Shortfall (99\%, 1-day) & 7{,}995{,}643 \\
Portfolio VaR (USD/EUR/GBP)      & 11{,}529{,}764 \\
Diversification benefit          & 9.6\% \\
\bottomrule
\end{tabular}
\end{table}

\noindent The illustrative engine demonstrates the mechanics cleanly but
rests on a critical assumption: that daily returns are normally distributed
with constant volatility.

% ---------- Section 5 ----------
\section{Calibration 2: Real-Data Engine}

\subsection{Data}

\begin{table}[h!]
\centering
\renewcommand{\arraystretch}{1.3}
\begin{tabular}{L{7cm} L{7cm}}
\toprule
\textbf{Item} & \textbf{Value} \\
\midrule
Source              & Bank of Ghana interbank FX rates \\
Instrument          & USD/GHS reference rate (mid) \\
Date range          & 2015-01-02 to 2025-12-31 \\
Observations        & 2{,}725 daily mid rates \\
Latest rate         & 10.4500 \\
Position value      & \ghs{209{,}000{,}000} \\
Raw file            & \texttt{data/usd\_ghs\_rates.csv} \\
Cleaned file        & \texttt{data/usd\_ghs\_clean.csv} \\
\bottomrule
\end{tabular}
\end{table}

\subsection{Empirical distribution diagnostics}

\begin{table}[h!]
\centering
\renewcommand{\arraystretch}{1.3}
\begin{tabular}{L{7cm} r}
\toprule
\textbf{Metric} & \textbf{Value} \\
\midrule
Realised daily volatility      & 0.7735\% \\
Realised annual volatility     & 12.2790\% \\
Kurtosis                       & \textbf{147.78} (normal = 0) \\
Skewness                       & $-1.84$ (normal = 0) \\
Worst single day               & $-15.04\%$ \\
Best single day                & $+13.34\%$ \\
\bottomrule
\end{tabular}
\end{table}

\noindent\textbf{Kurtosis of 147.78} is the central empirical fact. Under
a normal distribution, kurtosis is 0. A kurtosis of 147 implies that
extreme moves occur with vastly greater frequency than the normal
distribution allows. This alone predicts that a parametric VaR model
will fail.

\subsection{VaR estimates: parametric vs.\ historical}

\begin{table}[h!]
\centering
\renewcommand{\arraystretch}{1.3}
\begin{tabular}{L{4cm} r r}
\toprule
\textbf{Method} & \textbf{VaR (GHS)} & \textbf{ES (GHS)} \\
\midrule
Parametric (full sample)  & 3{,}760{,}816 & 4{,}308{,}633 \\
Historical (full sample)  & 3{,}951{,}268 & 9{,}229{,}126 \\
\midrule
\textbf{Ratio (Hist / Param)} & \textbf{1.05\(\times\)} & \textbf{2.14\(\times\)} \\
\bottomrule
\end{tabular}
\end{table}

\noindent The parametric and historical VaR estimates are close
(1.05\(\times\)), but the ES estimates diverge by 2.14\(\times\). This
is exactly what fat tails produce: the two methods agree on the body of
the distribution but the parametric method drastically understates the
tail.

\subsection{GARCH(1,1) conditional volatility}

The full-sample parametric VaR assumes a single constant volatility.
In reality, GHS volatility clusters. A GARCH(1,1) model with Student-\(t\)
errors was fitted to the 2{,}725 daily log returns:

\begin{table}[h!]
\centering
\renewcommand{\arraystretch}{1.3}
\begin{tabular}{L{4cm} r L{7cm}}
\toprule
\textbf{Parameter} & \textbf{Estimate} & \textbf{Interpretation} \\
\midrule
$\omega$     & $1.70 \times 10^{-4}$ & Baseline variance \\
$\alpha_1$   & 0.5287 & Reaction to yesterday's shock (large) \\
$\beta_1$    & 0.4713 & Volatility persistence (large) \\
$\alpha_1 + \beta_1$ & \textbf{1.000} & \textbf{IGARCH: shocks are permanent} \\
$\nu$ (t d.o.f.) & 2.83 & Extreme fat tails (normal is $\infty$) \\
\bottomrule
\end{tabular}
\end{table}

\noindent The sum $\alpha_1 + \beta_1 \approx 1.000$ indicates an
\textbf{integrated GARCH} process: volatility shocks to the cedi do not
decay --- they persist. This is consistent with a currency in structural
depreciation. The Student-\(t\) degrees of freedom of $\nu = 2.83$ confirms
that the conditional distribution is extremely fat-tailed.

\subsection{GARCH vs.\ parametric VaR}

\begin{table}[h!]
\centering
\renewcommand{\arraystretch}{1.3}
\begin{tabular}{L{6cm} r}
\toprule
\textbf{Measure} & \textbf{Value (GHS)} \\
\midrule
Full-sample daily volatility (constant) & 0.7735\% \\
GARCH next-day conditional $\sigma$     & \textbf{1.7909\%} \\
\midrule
Parametric VaR (99\%, 1-day)            & 3{,}760{,}816 \\
\textbf{GARCH VaR (99\%, 1-day)}        & \textbf{8{,}707{,}564} \\
\midrule
\textbf{GARCH / Parametric ratio}       & \textbf{2.32\(\times\)} \\
\bottomrule
\end{tabular}
\end{table}

\noindent The GARCH conditional VaR is \textbf{2.32\(\times\) the
parametric VaR}. The constant-volatility assumption materially understates
risk when the market is already in a stressed regime.

\subsection{Volatility regimes, 2015--2025}

\begin{table}[h!]
\centering
\renewcommand{\arraystretch}{1.2}
\begin{tabular}{r r r r}
\toprule
\textbf{Year} & \textbf{Daily vol} & \textbf{Annual vol} & \textbf{VaR 99\% 1-day (GHS)} \\
\midrule
2015 & 0.8535\% & 13.55\% & 4{,}149{,}900 \\
2016 & 0.1596\% &  2.53\% &   775{,}935 \\
2017 & 0.3070\% &  4.87\% & 1{,}492{,}586 \\
2018 & 0.1492\% &  2.37\% &   725{,}510 \\
2019 & 0.2546\% &  4.04\% & 1{,}237{,}767 \\
2020 & 0.1497\% &  2.38\% &   727{,}672 \\
2021 & 0.0552\% &  0.88\% &   268{,}605 \\
\textbf{2022} & \textbf{1.9397\%} & \textbf{30.79\%} & \textbf{9{,}430{,}984} \\
2023 & 0.9332\% & 14.81\% & 4{,}537{,}446 \\
2024 & 0.3003\% &  4.77\% & 1{,}460{,}052 \\
2025 & 0.9400\% & 14.92\% & 4{,}570{,}294 \\
\bottomrule
\end{tabular}
\end{table}

\noindent\textbf{35\(\times\) spread.} A VaR model calibrated on the
2021 regime (0.88\% annual) would have understated 2022 crisis risk by
a factor of 35. This is the primary justification for stress testing,
GARCH-based volatility models, and regime-conditional capital buffers.

\subsection{Two historical VaR windows, both reported}

The framework deliberately reports two historical VaR measures:

\begin{table}[h!]
\centering
\renewcommand{\arraystretch}{1.3}
\begin{tabular}{L{6cm} r}
\toprule
\textbf{Window} & \textbf{Historical VaR (GHS)} \\
\midrule
Full-sample (2{,}725 days)  & 3{,}951{,}268 \\
Rolling 250-day             & 9{,}581{,}018 \\
\midrule
\textbf{Recency uplift}     & \textbf{2.43\(\times\)} \\
\bottomrule
\end{tabular}
\end{table}

\noindent Both are correct. The full-sample figure is the long-run,
unconditional estimate used for capital planning and IMA regulatory
purposes. The rolling 250-day figure is the recent conditional estimate
used for daily limit monitoring. The recency uplift of \(2.43\times\)
reflects elevated 2025 volatility relative to the 10-year average. The
framework never presents either number in isolation.

\subsection{Kupiec backtest}

\begin{table}[h!]
\centering
\renewcommand{\arraystretch}{1.3}
\begin{tabular}{L{7cm} r}
\toprule
\textbf{Item} & \textbf{Value} \\
\midrule
Backtest window                & Rolling 250-day \\
Observations                   & 2{,}475 \\
Expected breaches              & 24.75 \\
\textbf{Observed breaches}     & \textbf{60} \\
Breach rate                    & 2.424\% \\
\(LR\) statistic               & \textbf{36.27} \\
Critical value \(\chi^2_{1,0.95}\) & 3.84 \\
\textbf{Verdict}               & \textbf{REJECTED} \\
\bottomrule
\end{tabular}
\end{table}

\noindent The parametric model is rejected at the 95\% test level.
The observed breach rate (2.42\%) is 2.42\(\times\) the expected rate
(1.00\%). This is not a coding error; it is the correct verdict for a
normal-distribution model applied to a kurtosis-147 return series.

% ---------- Section 6 ----------
\section{Comparison of the Two Calibrations}

\begin{table}[h!]
\centering
\renewcommand{\arraystretch}{1.3}
\begin{tabular}{L{5cm} r r}
\toprule
\textbf{Metric} & \textbf{Illustrative} & \textbf{Real Data} \\
\midrule
Position value                 & \ghs{250{,}000{,}000} & \ghs{209{,}000{,}000} \\
Assumed / realised daily vol   & 1.20\% & 0.7735\% \\
Parametric VaR (99\%, 1-day)   & 6{,}979{,}044 & 3{,}760{,}816 \\
Expected Shortfall             & 7{,}995{,}643 & 4{,}308{,}633 \\
Fat tails captured?            & No     & Yes (kurtosis 147.78) \\
Kupiec verdict                 & N/A    & \textbf{REJECTED} \\
Model survives?                & Unknown & \textbf{No} \\
\bottomrule
\end{tabular}
\end{table}

\noindent\textbf{What the comparison teaches.} The illustrative engine
shows \emph{how} to compute a VaR. The real-data engine shows
\emph{when the underlying assumptions break}. In a production risk
function at a Ghanaian bank, the parametric engine would be a first pass
only --- the operative measures would be historical simulation, Expected
Shortfall, GARCH, and stress testing.

% ---------- Section 7 ----------
\section{Stress Testing}

\subsection{Backward-Looking: 2022 Crisis Replay}

\begin{table}[h!]
\centering
\renewcommand{\arraystretch}{1.3}
\begin{tabular}{L{7cm} r}
\toprule
\textbf{Scenario} & \textbf{Value (GHS)} \\
\midrule
2022 daily volatility       & 1.9397\% \\
Stressed VaR (99\%, 1-day)  & 9{,}430{,}984 \\
Stressed ES (99\%, 1-day)   & 10{,}804{,}744 \\
Stress / full-sample ratio  & \textbf{2.51\(\times\)} \\
\bottomrule
\end{tabular}
\end{table}

\subsection{Forward-Looking: Overnight Depreciation Scenarios}

\begin{table}[h!]
\centering
\renewcommand{\arraystretch}{1.3}
\begin{tabular}{r r r r}
\toprule
\textbf{Shock} & \textbf{Shocked USD/GHS} & \textbf{Revaluation (GHS)} & \textbf{\% of Position} \\
\midrule
10\% & 11.4950 & 20{,}900{,}000 & 10.00\% \\
20\% & 12.5400 & 41{,}800{,}000 & 20.00\% \\
\textbf{30\%} & \textbf{13.5850} & \textbf{62{,}700{,}000} & \textbf{30.00\%} \\
\bottomrule
\end{tabular}
\end{table}

\noindent The 30\% scenario corresponds to the scale of the 2022 cedi
depreciation within a single day. For a bank funding a USD 20m long
position in GHS, the funding side would carry a GH₵62.7m loss. This is
the ALCO stress scenario recommended in Section~9.

% ---------- Section 8 ----------
\section{Ghana Regulatory Implications}

\subsection{Basel III / BoG capital charge}
\begin{equation}
\text{Capital Charge} = \max\!\left(\text{VaR}_{t-1},\;
m_c \cdot \text{VaR}_{avg,60}\right) + \text{SVaR}
\end{equation}
with \(m_c \geq 3\) and BoG's minimum CAR of 13\%.

\subsection{Backtesting consequences}

Under the Basel traffic-light framework:

\begin{table}[h!]
\centering
\renewcommand{\arraystretch}{1.3}
\begin{tabular}{L{3cm} L{5cm} L{5cm}}
\toprule
\textbf{Zone} & \textbf{Breach range} & \textbf{Capital consequence} \\
\midrule
\textcolor{ghgreen}{\textbf{Green}}  & 0--4   & Multiplier stays at 3.0 \\
\textcolor{ghgold}{\textbf{Yellow}}  & 5--9   & Multiplier rises to 3.4--3.85 \\
\textcolor{ghred}{\textbf{Red}}      & 10+    & Model rejected; BoG intervention \\
\bottomrule
\end{tabular}
\end{table}

\noindent A real-data breach count of 60 over 2{,}475 rolling windows
corresponds to a rolling 250-day breach count of \(\approx 6\)---placing
the parametric model in the \textbf{yellow zone} and increasing the
regulatory capital multiplier.

% ---------- Section 9 ----------
\section{Recommendations}

\begin{enumerate}[leftmargin=1.6em]
    \item Migrate from parametric to \textbf{historical simulation} as
          the primary daily VaR measure for the FX book.
    \item Report \textbf{Expected Shortfall} alongside VaR; the empirical
          ES is 2.14\(\times\) the parametric ES.
    \item Adopt a \textbf{GARCH(1,1)} volatility model to capture the
          IGARCH persistence observed in the data.
    \item Extend stress testing to include a \textbf{30\% overnight GHS
          depreciation} scenario for ALCO.
    \item Present this dual-calibration framework to the model validation
          committee for review.
    \item Publish the analysis openly (\texttt{\reporef}) as a
          contribution to Ghanaian quantitative finance.
\end{enumerate}

% ---------- Section 10 ----------
\section{Key Assumptions and Limitations}

\begin{enumerate}[leftmargin=1.6em]
    \item \textbf{Normality.} Assumed in the parametric engine; falsified
          by the real-data kurtosis of 147.78.
    \item \textbf{Constant volatility.} The parametric model assumes it;
          GARCH shows $\alpha_1 + \beta_1 = 1.000$ (integrated).
    \item \textbf{Linear positions.} No optionality or convexity.
    \item \textbf{Square-root-of-time scaling.} Assumes i.i.d.\ returns;
          GHS autocorrelation mildly violates this.
    \item \textbf{Liquidity.} Model assumes exit at market prices; GHS
          interbank liquidity can disappear under stress.
    \item \textbf{Illustrative bank.} GCIB is fictional; only the
          methodology and the data are real.
\end{enumerate}

% ---------- Appendix ----------
\appendix
\section{Companion Repository}

All source code, data, and outputs are maintained at:

\begin{center}
\texttt{\reporef}
\end{center}

\subsection*{Repository contents}

\begin{table}[h!]
\centering
\renewcommand{\arraystretch}{1.3}
\begin{tabular}{L{6cm} L{8cm}}
\toprule
\textbf{File / Directory} & \textbf{Purpose} \\
\midrule
\texttt{src/ghana\_fx\_var.py}          & Illustrative VaR engine \\
\texttt{src/load\_bog\_data.py}         & BoG CSV loader and diagnostics \\
\texttt{src/var\_real\_data.py}         & Real-data VaR engine \\
\texttt{src/garch\_var.py}              & GARCH(1,1) conditional VaR \\
\texttt{src/daily\_risk\_report.py}     & Daily risk report (historical primary) \\
\texttt{data/usd\_ghs\_rates.csv}       & Raw BoG interbank rates (2015--2025) \\
\texttt{data/usd\_ghs\_clean.csv}       & Cleaned series with log returns \\
\texttt{docs/Model\_Validation\_Memo.tex} & Model validation memo \\
\texttt{examples/}                      & Terminal outputs for all engines \\
\bottomrule
\end{tabular}
\end{table}

\subsection*{Author}

\textbf{\authorname}\\
\authortitle\\
\authordesk\\
\texttt{\authoremail}

% ---------- Sign-off ----------
\vspace{0.8cm}
\noindent\rule{\textwidth}{0.4pt}

\begin{table}[h!]
\centering
\renewcommand{\arraystretch}{1.6}
\begin{tabular}{L{5cm} L{4cm} L{5cm}}
\toprule
\textbf{Role} & \textbf{Name} & \textbf{Signature / Date} \\
\midrule
Prepared by (Quant Analyst)     & \textbf{\authorname} &  \\
Reviewed by (Head of Market Risk) &  &  \\
Approved by (CRO)               &  &  \\
Filed with Bank of Ghana        &  &  \\
\bottomrule
\end{tabular}
\end{table}

\vfill
\begin{center}
\small\itshape
--- End of Report ---\\
Ghana Commercial \& Investment Bank \;|\; Quantitative Risk \& Financial Engineering\\
Prepared under Bank of Ghana Risk Management Directive\\[0.3cm]
\textbf{\authorname} \;|\; \authortitle
\end{center}

\end{document}